\documentclass[12pt,a4paper]{article}

% --- 1. PACKAGES --------------------------------------------------------------
\usepackage[utf8]{inputenc}
\usepackage{mathtools}
\usepackage[T1]{fontenc}
\usepackage{amsmath, amssymb}
\usepackage{geometry}
\usepackage{graphicx}
\usepackage{xcolor}
\usepackage{listings}
\usepackage{hyperref}
\usepackage{booktabs}
\usepackage{enumitem}
\usepackage{tcolorbox}
\usepackage{mdframed}
\usepackage{fancyhdr}
\usepackage{titlesec}

% Font Settings
\usepackage{newtxtext,newtxmath}
\usepackage{courier}

\geometry{margin=2.5cm}

% --- 2. WARNA TEMA ------------------------------------------------------------
\definecolor{codebg}{RGB}{245,247,250}
\definecolor{codeframe}{RGB}{180,195,215}
\definecolor{keywordcolor}{RGB}{0,112,192}
\definecolor{commentcolor}{RGB}{100,160,100}
\definecolor{stringcolor}{RGB}{180,70,50}
\definecolor{notebg}{RGB}{255,249,230}
\definecolor{noteframe}{RGB}{230,180,50}
\definecolor{tipbg}{RGB}{230,245,235}
\definecolor{tipframe}{RGB}{50,160,90}
\definecolor{sectioncolor}{RGB}{30,80,150}

% --- 3. STYLE SECTION & TOC ---------------------------------------------------
\titleformat{\section}{\Large\bfseries\color{sectioncolor}}{}{0em}{\thesection\quad}[\titlerule]
\titleformat{\subsection}{\large\bfseries\color{sectioncolor!80}}{}{0em}{\thesubsection\quad}
\titleformat{\subsubsection}{\normalsize\bfseries\color{sectioncolor!60}}{}{0em}{\thesubsubsection\quad}

\usepackage{tocloft}
\renewcommand{\cftsecleader}{\cftdotfill{\cftdotsep}}
\renewcommand{\cftsubsecleader}{\cftdotfill{\cftdotsep}}
\renewcommand{\cftsubsubsecleader}{\cftdotfill{\cftdotsep}}
\renewcommand{\cftsecfont}{\bfseries\color{sectioncolor}}
\setcounter{tocdepth}{3}

% --- 4. LISTINGS (CODE STYLE) -------------------------------------------------
\lstdefinestyle{customstyle}{
  basicstyle=\ttfamily\small,
  backgroundcolor=\color{codebg},
  frame=single,
  rulecolor=\color{codeframe},
  keywordstyle=\bfseries\color{keywordcolor},
  commentstyle=\itshape\color{commentcolor},
  stringstyle=\color{stringcolor},
  numbers=left,
  numberstyle=\tiny\color{gray},
  breaklines=true,
  tabsize=4,
  showstringspaces=false
}
\lstset{style=customstyle}

% --- 5. KOTAK CATATAN ---------------------------------------------------------
\newmdenv[
  backgroundcolor=notebg, linecolor=noteframe, linewidth=2pt,
  leftline=true, topline=false, rightline=false, bottomline=false,
  innerleftmargin=10pt, innertopmargin=6pt, innerbottommargin=6pt, skipabove=8pt
]{notebox}

\newmdenv[
  backgroundcolor=tipbg, linecolor=tipframe, linewidth=2pt,
  leftline=true, topline=false, rightline=false, bottomline=false,
  innerleftmargin=10pt, innertopmargin=6pt, innerbottommargin=6pt, skipabove=8pt
]{tipbox}

% --- 6. HEADER & FOOTER -------------------------------------------------------
\pagestyle{fancy}
\fancyhf{}
\fancyhead[L]{\small\textit{Pengolahan Citra Digital -- Bab 4}}
\fancyhead[R]{\small\textit{noteify -- Pemrosesan Domain Frekuensi}}
\fancyfoot[C]{\thepage}
\renewcommand{\headrulewidth}{0.4pt}

% --- 7. DATA COVER ------------------------------------------------------------
\title{
  {\LARGE\bfseries Pengolahan Citra Digital}\\[0.5em]
  {\large Bab 4: Pemrosesan Citra dalam Domain Frekuensi}\\[0.8em]
  {\normalsize Memahami Transformasi Fourier, teorema sampling, dan teknik pemfilteran
   berbasis frekuensi sebagai pendekatan alternatif yang powerful dalam pengolahan citra.}
}
\author{}
\date{}

% =============================================================================
\begin{document}

\maketitle
\thispagestyle{fancy}
\tableofcontents
\newpage

% =============================================================================
\section{Pendahuluan}
% =============================================================================

Pada Bab 3 kita memproses citra langsung pada nilai piksel (\textit{domain spasial}).
Bab 4 memperkenalkan pendekatan yang sama sekali berbeda: memindahkan citra ke
\textbf{domain frekuensi}, melakukan operasi di sana, lalu memindahkannya kembali.

Analogi yang tepat: bayangkan sebuah lagu. Kita bisa mendengarnya sebagai gelombang
tekanan udara yang berubah terhadap waktu (\textit{domain waktu}), atau kita bisa melihat
\textit{lembar musik}-nya yang menunjukkan nada dan amplitudo setiap frekuensi
(\textit{domain frekuensi}). Keduanya merepresentasikan informasi yang \textit{sama},
hanya caranya berbeda.

Dalam konteks citra:
\begin{itemize}
  \item \textbf{Frekuensi rendah} merepresentasikan variasi lambat -- latar belakang,
        gradien halus, dan area dengan intensitas seragam.
  \item \textbf{Frekuensi tinggi} merepresentasikan variasi cepat -- tepi (\textit{edges}),
        noise, dan detail halus.
\end{itemize}

Keunggulan domain frekuensi: konvolusi dengan kernel besar yang mahal di domain spasial
menjadi sekadar perkalian elemen-per-elemen di domain frekuensi, jauh lebih efisien.

% =============================================================================
\section{Intuisi Dasar: Ide Besar Fourier (Bagian 4.1 \& 4.2)}
% =============================================================================

\subsection{Teorema Fourier}

Jean Baptiste Joseph Fourier (1807) mengemukakan ide yang pada masanya dianggap
radikal: \textit{setiap fungsi periodik dapat direpresentasikan sebagai penjumlahan
tak terhingga dari sines dan cosines dengan frekuensi berbeda}, masing-masing
dikalikan dengan koefisien tertentu. Ide ini kemudian digeneralisasikan untuk
fungsi non-periodik melalui \textbf{Transformasi Fourier}.

\subsection{Pasangan Transformasi Fourier 1-D Kontinu}

Untuk fungsi kontinu $f(t)$, \textbf{Transformasi Fourier Maju} didefinisikan sebagai:

\begin{equation}
  F(\mu) = \int_{-\infty}^{\infty} f(t)\, e^{-j2\pi\mu t}\, dt
  \label{eq:ft_forward}
\end{equation}

Dan \textbf{Transformasi Fourier Balik} (Inverse FT):

\begin{equation}
  f(t) = \int_{-\infty}^{\infty} F(\mu)\, e^{j2\pi\mu t}\, d\mu
  \label{eq:ft_inverse}
\end{equation}

di mana $\mu$ mewakili frekuensi (dalam Hz jika $t$ dalam detik) dan
$j = \sqrt{-1}$. Berdasarkan rumus Euler: $e^{j\theta} = \cos\theta + j\sin\theta$,
sehingga setiap komponen Fourier pada dasarnya adalah gelombang sinus/cosinus.

\begin{notebox}
\textbf{Intuisi Persamaan~\eqref{eq:ft_forward}:} Faktor $e^{-j2\pi\mu t}$
berfungsi sebagai ``detektor'' frekuensi. Mengalikan $f(t)$ dengan detektor
frekuensi $\mu$ lalu mengintegrasikan akan mengekstrak seberapa kuat komponen
frekuensi $\mu$ terdapat di dalam $f(t)$.
\end{notebox}

\subsection{Representasi Kompleks: Spektrum, Fase, dan Daya}

Karena $F(\mu)$ adalah bilangan kompleks, ia dapat ditulis sebagai:

\begin{equation}
  F(\mu) = R(\mu) + j\,I(\mu)
  \label{eq:ft_complex}
\end{equation}

\textbf{Spektrum Amplitudo} (yang umum divisualisasikan):

\begin{equation}
  |F(\mu)| = \sqrt{R^2(\mu) + I^2(\mu)}
  \label{eq:spectrum_1d}
\end{equation}

\textbf{Sudut Fase} (menyimpan informasi lokasi dan struktur):

\begin{equation}
  \phi(\mu) = \arctan\!\left(\frac{I(\mu)}{R(\mu)}\right)
  \label{eq:phase_1d}
\end{equation}

\textbf{Spektrum Daya} (\textit{Power Spectrum}):

\begin{equation}
  P(\mu) = |F(\mu)|^2 = R^2(\mu) + I^2(\mu)
  \label{eq:power_1d}
\end{equation}

\subsection{Sifat-Sifat Penting Transformasi Fourier}

\begin{center}
\renewcommand{\arraystretch}{1.4}
\begin{tabular}{lll}
\toprule
\textbf{Sifat} & \textbf{Domain Spasial/Waktu} & \textbf{Domain Frekuensi} \\
\midrule
Linearitas     & $af(t) + bg(t)$               & $aF(\mu) + bG(\mu)$ \\
Pergeseran     & $f(t - t_0)$                  & $e^{-j2\pi\mu t_0}\,F(\mu)$ \\
Konvolusi      & $f(t) * h(t)$                 & $F(\mu) \cdot H(\mu)$ \\
Korelasi       & $f(t) \star g(t)$             & $F^*(\mu) \cdot G(\mu)$ \\
Simetri        & $f(t)$ riil                   & $F(-\mu) = F^*(\mu)$ \\
\bottomrule
\end{tabular}
\end{center}

\begin{tipbox}
\textbf{Sifat Konvolusi adalah kunci utama bab ini:} konvolusi di domain spasial
setara dengan perkalian di domain frekuensi (lihat Persamaan~\eqref{eq:conv_theorem}).
Ini mengubah operasi yang mahal menjadi perkalian sederhana setelah FFT.
\end{tipbox}

% =============================================================================
\section{Sampling dan Teorema Sampling (Bagian 4.3)}
% =============================================================================

\subsection{Kebutuhan Sampling}

Sebelum citra dapat diolah komputer, sinyal kontinu harus di-\textit{sampling} --
diambil nilainya pada titik-titik diskrit berjarak $\Delta T$. Pertanyaan kritis:
\textbf{seberapa rapat sampling agar tidak ada informasi yang hilang?}

\subsection{Teorema Sampling Nyquist-Shannon}

Sebuah fungsi yang dibatasi pita frekuensinya (\textit{band-limited}) dengan
frekuensi tertinggi $\mu_{\max}$ dapat dipulihkan \textit{secara sempurna} dari
sampel-sampelnya jika dan hanya jika laju sampling memenuhi:

\begin{equation}
  \frac{1}{\Delta T} > 2\,\mu_{\max}
  \label{eq:nyquist}
\end{equation}

Nilai $2\,\mu_{\max}$ disebut \textbf{Laju Nyquist}. Interval sampling
$\Delta T = \tfrac{1}{2\mu_{\max}}$ disebut \textbf{Interval Nyquist}.

\subsection{Aliasing}

Jika Persamaan~\eqref{eq:nyquist} dilanggar (sampling terlalu jarang), terjadi
\textbf{aliasing}: komponen frekuensi tinggi ``berpura-pura'' menjadi frekuensi
rendah yang salah. Secara matematis, spektrum sinyal yang di-sampling bersifat
periodik, dan jika laju sampling tidak cukup, periode-periode tersebut akan
\textit{saling tumpang tindih} (\textit{spectral overlap}).

\begin{notebox}
Lihat ilustrasi visual aliasing pada gelombang sinus di \textbf{[FIGURE 4.11]}.
\end{notebox}

\begin{tipbox}
\textbf{Contoh Aliasing di Dunia Nyata:}
\begin{itemize}
  \item Roda berputar di film yang tampak berputar mundur (\textit{wagon-wheel effect}).
  \item Garis diagonal pada layar yang tampak bergerigi (\textit{jaggies}).
  \item Pola Moire pada foto kain bermotif halus.
\end{itemize}
Solusi umum: terapkan \textit{anti-aliasing filter} (lowpass) sebelum sampling.
\end{tipbox}

\subsection{Rekonstruksi Sinyal: Interpolasi Sinc}

Sinyal kontinu dapat direkonstruksi sempurna dari sampelnya menggunakan:

\begin{equation}
  f(t) = \sum_{n=-\infty}^{\infty} f(n\Delta T)\;
         \mathrm{sinc}\!\left(\frac{t - n\Delta T}{\Delta T}\right)
  \label{eq:reconstruction}
\end{equation}

di mana $\mathrm{sinc}(x) = \dfrac{\sin(\pi x)}{\pi x}$.

% =============================================================================
\section{Discrete Fourier Transform 2-D (Bagian 4.5 \& 4.6)}
% =============================================================================

\subsection{DFT 1-D Diskrit}

Untuk sinyal diskrit $f(x)$ dengan $N$ sampel, DFT 1-D didefinisikan sebagai:

\begin{equation}
  F(u) = \sum_{x=0}^{N-1} f(x)\, e^{-j2\pi ux/N},
  \qquad u = 0, 1, \ldots, N-1
  \label{eq:dft_1d}
\end{equation}

Dan inversnya (IDFT 1-D):

\begin{equation}
  f(x) = \frac{1}{N} \sum_{u=0}^{N-1} F(u)\, e^{j2\pi ux/N},
  \qquad x = 0, 1, \ldots, N-1
  \label{eq:idft_1d}
\end{equation}

\subsection{DFT 2-D untuk Citra $M \times N$}

\begin{equation}
  F(u,v) = \sum_{x=0}^{M-1} \sum_{y=0}^{N-1}
            f(x,y)\, e^{-j2\pi\!\left(\frac{ux}{M} + \frac{vy}{N}\right)}
  \label{eq:dft_2d}
\end{equation}

Dan inversnya (IDFT 2-D):

\begin{equation}
  f(x,y) = \frac{1}{MN} \sum_{u=0}^{M-1} \sum_{v=0}^{N-1}
            F(u,v)\, e^{\,j2\pi\!\left(\frac{ux}{M} + \frac{vy}{N}\right)}
  \label{eq:idft_2d}
\end{equation}

\subsection{Spektrum Amplitudo 2-D}

Karena $F(u,v)$ adalah bilangan kompleks, spektrum amplitudo yang divisualisasikan:

\begin{equation}
  |F(u,v)| = \sqrt{R^2(u,v) + I^2(u,v)}
  \label{eq:spectrum_2d}
\end{equation}

Karena rentang nilainya sangat besar, biasanya ditampilkan dalam skala logaritmik:

\begin{equation}
  D(u,v) = c \cdot \log\!\bigl(1 + |F(u,v)|\bigr)
  \label{eq:spectrum_log}
\end{equation}

\subsection{Sudut Fase 2-D}

\begin{equation}
  \phi(u,v) = \arctan\!\left(\frac{I(u,v)}{R(u,v)}\right)
  \label{eq:phase_2d}
\end{equation}

\begin{notebox}
\textbf{Fase vs.\ Amplitudo:} Jika kita merekonstruksi citra hanya menggunakan
informasi fase (amplitudo diset konstan), bentuk objek masih dapat dikenali.
Sebaliknya, jika hanya menggunakan amplitudo (fase diset nol), citra menjadi
kacau. Ini menunjukkan bahwa \textbf{fase menyimpan informasi struktural} citra.
Lihat \textbf{[FIGURE 4.26]}.
\end{notebox}

\subsection{Periodisitas dan Centering}

DFT bersifat periodik dengan periode $M$ dan $N$. Akibatnya, frekuensi rendah
(komponen DC) berada di sudut-sudut spektrum, bukan di tengah. Untuk memudahkan
visualisasi dan pemfilteran, spektrum digeser agar frekuensi rendah berada di
tengah dengan mengalikan citra sebelum DFT dengan:

\begin{equation}
  f_c(x,y) = f(x,y) \cdot (-1)^{x+y}
  \label{eq:centering}
\end{equation}

Setelah pemrosesan selesai, kalikan kembali dengan $(-1)^{x+y}$ untuk
mengembalikan ke posisi semula.

\subsection{Fast Fourier Transform (FFT)}

DFT naif memerlukan $O(N^2)$ operasi untuk sinyal 1-D berukuran $N$.
\textbf{FFT} (algoritma Cooley-Tukey, 1965) mereduksinya menjadi $O(N \log N)$.
Untuk citra $N \times N$, DFT 2-D naif memerlukan $O(N^4)$, sementara FFT 2-D
hanya $O(N^2 \log N)$ -- perbedaan yang sangat signifikan pada citra besar.

\begin{center}
\renewcommand{\arraystretch}{1.3}
\begin{tabular}{lrr}
\toprule
\textbf{Ukuran Citra} & \textbf{DFT Naif (operasi)} & \textbf{FFT (operasi)} \\
\midrule
$256 \times 256$   & $\approx 4 \times 10^{9}$  & $\approx 1 \times 10^{6}$ \\
$512 \times 512$   & $\approx 7 \times 10^{10}$ & $\approx 4 \times 10^{6}$ \\
$1024 \times 1024$ & $\approx 1 \times 10^{12}$ & $\approx 2 \times 10^{7}$ \\
\bottomrule
\end{tabular}
\end{center}

% =============================================================================
\section{Alur Kerja Pemfilteran Domain Frekuensi (Bagian 4.7)}
% =============================================================================

\subsection{Teorema Konvolusi sebagai Landasan}

Sifat konvolusi dari Transformasi Fourier menyatakan:

\begin{equation}
  f(x,y) * h(x,y) \;\xleftrightarrow{\;\mathcal{F}\;}\; F(u,v) \cdot H(u,v)
  \label{eq:conv_theorem}
\end{equation}

Persamaan~\eqref{eq:conv_theorem} adalah motivasi utama pemfilteran domain
frekuensi: konvolusi spasial yang kompleks diganti dengan perkalian sederhana
setelah FFT.

\subsection{Langkah-Langkah Lengkap Pemfilteran}

Diberikan citra masukan $f(x,y)$ berukuran $M \times N$:

\begin{enumerate}
  \item \textbf{Padding.} Perluas citra menjadi $P \times Q$ dengan menambah
        nol (\textit{zero-padding}):
        \begin{equation}
          P \geq 2M - 1, \qquad Q \geq 2N - 1
          \label{eq:padding_size}
        \end{equation}
        Biasanya dibulatkan ke pangkat 2 terdekat untuk efisiensi FFT.

  \item \textbf{Centering.} Kalikan dengan $(-1)^{x+y}$
        (Persamaan~\eqref{eq:centering}).

  \item \textbf{DFT.} Hitung:
        \begin{equation}
          F_p(u,v) = \mathcal{F}\!\left\{f_p(x,y)\right\}
          \label{eq:step_dft}
        \end{equation}

  \item \textbf{Filter.} Kalikan dengan fungsi filter:
        \begin{equation}
          G(u,v) = H(u,v) \cdot F_p(u,v)
          \label{eq:filtering_step}
        \end{equation}

  \item \textbf{IDFT.} Hitung:
        \begin{equation}
          g_p(x,y) = \mathcal{F}^{-1}\!\left\{G(u,v)\right\}
          \label{eq:step_idft}
        \end{equation}

  \item \textbf{Post-processing.} Ambil bagian riil, kembalikan centering,
        dan potong ke ukuran asli:
        \begin{equation}
          g(x,y) = \operatorname{Re}\!\left[g_p(x,y) \cdot (-1)^{x+y}\right]
                   \Big|_{\substack{0 \leq x \leq M-1 \\ 0 \leq y \leq N-1}}
          \label{eq:postprocess}
        \end{equation}
\end{enumerate}

\begin{notebox}
\textbf{Mengapa Perlu Padding?} Tanpa padding, DFT menghasilkan \textit{circular
convolution} (konvolusi siklik): piksel dari sisi kanan ``bocor'' ke sisi kiri.
Padding memastikan konvolusi siklik bersifat identik dengan konvolusi linear pada
region yang kita inginkan, menghindari \textit{wraparound error}.
\end{notebox}

\subsection{Contoh Kode: Alur Kerja Lengkap FFT Filtering}

\begin{lstlisting}[language=Python, caption={Alur kerja pemfilteran domain frekuensi (6 langkah)}]
import numpy as np
from PIL import Image
import matplotlib.pyplot as plt

def freq_domain_filter(img, H):
    """
    Terapkan filter H(u,v) ke citra menggunakan domain frekuensi.
    img : numpy array 2D (grayscale)
    H   : numpy array 2D, ukuran P x Q (setelah padding)
    """
    M, N = img.shape
    P, Q = H.shape

    # Langkah 1: Padding
    fp = np.zeros((P, Q))
    fp[:M, :N] = img

    # Langkah 2: Centering -- (-1)^(x+y)
    x, y = np.arange(P), np.arange(Q)
    xx, yy = np.meshgrid(x, y, indexing='ij')
    center_mask = (-1.0) ** (xx + yy)
    fp_c = fp * center_mask

    # Langkah 3: DFT
    F = np.fft.fft2(fp_c)

    # Langkah 4: Filter G(u,v) = H(u,v) * F(u,v)
    G = H * F

    # Langkah 5: IDFT
    gp = np.fft.ifft2(G)

    # Langkah 6: Post-processing
    g = np.real(gp) * center_mask   # balik centering
    g = g[:M, :N]                   # potong ke ukuran asli

    return np.clip(g, 0, 255)


# --- Demo dengan Ideal Lowpass Filter ---
img = np.array(Image.open("image.jpg").convert("L"), dtype=np.float64)
M, N = img.shape
P, Q = 2 * M, 2 * N

# Buat filter Ideal Lowpass (centered)
D0 = 50
u = np.fft.fftfreq(P) * P
v = np.fft.fftfreq(Q) * Q
uu, vv = np.meshgrid(u, v, indexing='ij')
D = np.sqrt(uu**2 + vv**2)
H_ilpf = (D <= D0).astype(np.float64)

result = freq_domain_filter(img, H_ilpf)

fig, axes = plt.subplots(1, 2, figsize=(10, 5))
axes[0].imshow(img.astype(np.uint8), cmap='gray')
axes[0].set_title('Citra Asli')
axes[1].imshow(result.astype(np.uint8), cmap='gray')
axes[1].set_title('Setelah ILPF (D0=50)')
for ax in axes:
    ax.axis('off')
plt.tight_layout()
plt.show()
\end{lstlisting}

% =============================================================================
\section{Jenis-Jenis Filter (Bagian 4.8 \& 4.9)}
% =============================================================================

\subsection{Jarak dari Pusat Spektrum}

Dasar dari semua filter frekuensi radial adalah jarak setiap titik $(u,v)$ dari
pusat (DC component):

\begin{equation}
  D(u,v) = \sqrt{u^2 + v^2}
  \label{eq:distance}
\end{equation}

Seluruh filter lowpass dan highpass isotropik didefinisikan berdasarkan
Persamaan~\eqref{eq:distance} dan nilai cutoff $D_0$.

\subsection{Hubungan Lowpass dan Highpass}

Setiap filter highpass dapat diturunkan langsung dari filter lowpass:

\begin{equation}
  H_{HP}(u,v) = 1 - H_{LP}(u,v)
  \label{eq:hp_from_lp}
\end{equation}

\subsection{Filter Lowpass (Smoothing)}

\subsubsection{Ideal Lowpass Filter (ILPF)}

\begin{equation}
  H_{\mathrm{ILPF}}(u,v) =
  \begin{cases}
    1 & \text{jika } D(u,v) \leq D_0 \\
    0 & \text{jika } D(u,v) > D_0
  \end{cases}
  \label{eq:ilpf}
\end{equation}

Transisi yang sangat tajam di domain frekuensi (Persamaan~\eqref{eq:ilpf})
menghasilkan osilasi (\textit{ringing} / \textit{Gibbs phenomenon}) di domain
spasial -- berupa lingkaran-lingkaran di sekitar tepi objek.

\begin{notebox}
Lihat efek \textit{ringing} dari ILPF pada \textbf{[FIGURE 4.41]}.
\end{notebox}

\subsubsection{Butterworth Lowpass Filter (BLPF)}

\begin{equation}
  H_{\mathrm{BLPF}}(u,v) = \frac{1}{1 + \left[\dfrac{D(u,v)}{D_0}\right]^{2n}}
  \label{eq:blpf}
\end{equation}

di mana $n$ adalah orde filter. Semakin besar $n$, semakin tajam transisi
dan semakin mirip ILPF, namun ringing juga semakin muncul.

\begin{tipbox}
\textbf{Panduan Pemilihan Orde:} $n = 1$ tidak menghasilkan ringing;
$n = 2$ merupakan kompromi yang baik; $n \geq 5$ mulai menunjukkan ringing.
\end{tipbox}

\subsubsection{Gaussian Lowpass Filter (GLPF)}

\begin{equation}
  H_{\mathrm{GLPF}}(u,v) = e^{-D^2(u,v)\,/\,(2D_0^2)}
  \label{eq:glpf}
\end{equation}

\textbf{Keunggulan unik:} Transformasi Fourier dari fungsi Gaussian adalah
Gaussian itu sendiri. Ini berarti GLPF di domain frekuensi berkorespondensi
dengan Gaussian filter di domain spasial, dan Gaussian \textbf{tidak menghasilkan
ringing sama sekali}.

\begin{notebox}
Perbandingan ketiga filter lowpass dapat dilihat pada \textbf{[FIGURE 4.44]}.
\end{notebox}

\subsection{Filter Highpass (Sharpening)}

Menggunakan Persamaan~\eqref{eq:hp_from_lp}, ketiga filter highpass:

\begin{equation}
  H_{\mathrm{IHPF}}(u,v) =
  \begin{cases}
    0 & \text{jika } D(u,v) \leq D_0 \\
    1 & \text{jika } D(u,v) > D_0
  \end{cases}
  \label{eq:ihpf}
\end{equation}

\begin{equation}
  H_{\mathrm{BHPF}}(u,v) = \frac{1}{1 + \left[\dfrac{D_0}{D(u,v)}\right]^{2n}}
  \label{eq:bhpf}
\end{equation}

\begin{equation}
  H_{\mathrm{GHPF}}(u,v) = 1 - e^{-D^2(u,v)\,/\,(2D_0^2)}
  \label{eq:ghpf}
\end{equation}

\subsection{High-Frequency Emphasis (HFE) Filter}

Filter highpass murni menghilangkan seluruh komponen DC sehingga citra tampak
``datar'' tanpa tonalitas. \textbf{HFE Filter} mengatasi ini dengan
mempertahankan sebagian komponen frekuensi rendah:

\begin{equation}
  H_{\mathrm{HFE}}(u,v) = a + b \cdot H_{HP}(u,v), \quad a \geq 0,\; b > a
  \label{eq:hfe}
\end{equation}

di mana $a$ adalah offset (biasanya $0 < a < 1$) dan $b$ adalah faktor
penguatan frekuensi tinggi.

\begin{notebox}
Jika dikombinasikan dengan Histogram Equalization setelah filtering, hasilnya
sangat efektif untuk meningkatkan ketajaman dan kontras secara bersamaan.
Lihat \textbf{[FIGURE 4.57]}.
\end{notebox}

\subsection{Contoh Kode: Perbandingan Filter Lowpass dan Visualisasi Spektrum}

\begin{lstlisting}[language=Python, caption={Perbandingan ILPF, BLPF, GLPF dan visualisasi spektrum}]
import numpy as np
from PIL import Image
import matplotlib.pyplot as plt

def make_distance_matrix(P, Q):
    u = np.fft.fftfreq(P) * P
    v = np.fft.fftfreq(Q) * Q
    uu, vv = np.meshgrid(u, v, indexing='ij')
    D = np.sqrt(uu**2 + vv**2)
    D[D == 0] = 1e-10
    return D

def make_lowpass_filters(P, Q, D0, n=2):
    D = make_distance_matrix(P, Q)
    H_ideal  = (D <= D0).astype(float)
    H_butter = 1.0 / (1.0 + (D / D0) ** (2 * n))
    H_gauss  = np.exp(-(D**2) / (2 * D0**2))
    return H_ideal, H_butter, H_gauss

def apply_filter(img, H):
    M, N = img.shape
    P, Q = H.shape
    fp = np.zeros((P, Q))
    fp[:M, :N] = img
    x, y = np.arange(P), np.arange(Q)
    xx, yy = np.meshgrid(x, y, indexing='ij')
    mask = (-1.0) ** (xx + yy)
    F = np.fft.fft2(fp * mask)
    G = H * F
    g = np.real(np.fft.ifft2(G)) * mask
    return np.clip(g[:M, :N], 0, 255)

img = np.array(Image.open("image.jpg").convert("L"), dtype=np.float64)
M, N = img.shape
P, Q = 2 * M, 2 * N
D0 = 30

# Spektrum Fourier
x, y = np.arange(M), np.arange(N)
xx, yy = np.meshgrid(x, y, indexing='ij')
mask = (-1.0) ** (xx + yy)
F = np.fft.fft2(img * mask)
spectrum = np.log(1 + np.abs(F))
spectrum = 255 * spectrum / spectrum.max()

# Filter
H_ideal, H_butter, H_gauss = make_lowpass_filters(P, Q, D0, n=2)

fig, axes = plt.subplots(2, 3, figsize=(15, 8))

axes[0, 0].imshow(img.astype(np.uint8), cmap='gray')
axes[0, 0].set_title('Citra Asli')
axes[0, 1].imshow(spectrum.astype(np.uint8), cmap='gray')
axes[0, 1].set_title('Spektrum Fourier (log)')
axes[0, 2].axis('off')

axes[1, 0].imshow(apply_filter(img, H_ideal).astype(np.uint8), cmap='gray')
axes[1, 0].set_title('ILPF (D0=30)')
axes[1, 1].imshow(apply_filter(img, H_butter).astype(np.uint8), cmap='gray')
axes[1, 1].set_title('BLPF n=2 (D0=30)')
axes[1, 2].imshow(apply_filter(img, H_gauss).astype(np.uint8), cmap='gray')
axes[1, 2].set_title('GLPF (D0=30)')

for ax in axes.flatten():
    ax.axis('off')
plt.tight_layout()
plt.show()
\end{lstlisting}

% =============================================================================
\section{Homomorphic Filtering (Bagian 4.9)}
% =============================================================================

\subsection{Model Pencahayaan-Reflektansi}

Sebuah citra $f(x,y)$ dimodelkan sebagai perkalian dua komponen:

\begin{equation}
  f(x,y) = i(x,y) \cdot r(x,y)
  \label{eq:illumination_model}
\end{equation}

di mana:
\begin{itemize}
  \item $i(x,y)$ adalah \textbf{iluminasi} -- seberapa banyak cahaya mengenai
        objek. Bervariasi \textit{lambat} secara spasial $\Rightarrow$ frekuensi
        rendah. Rentang: $0 < i(x,y) < \infty$.
  \item $r(x,y)$ adalah \textbf{reflektansi} -- seberapa banyak cahaya
        dipantulkan objek. Bervariasi \textit{cepat} $\Rightarrow$ frekuensi
        tinggi. Rentang: $0 < r(x,y) < 1$.
\end{itemize}

\textbf{Masalah utama:} Transformasi Fourier dari perkalian bukan perkalian dari
Transformasi Fourier:
\[
  \mathcal{F}\{i \cdot r\} \neq \mathcal{F}\{i\} \cdot \mathcal{F}\{r\}
\]
Sehingga kita tidak bisa memisahkan $i$ dan $r$ secara langsung di domain
frekuensi.

\subsection{Solusi: Linearisasi dengan Logaritma}

Dengan mengambil logaritma natural, perkalian berubah menjadi penjumlahan:

\begin{equation}
  z(x,y) = \ln f(x,y) = \ln i(x,y) + \ln r(x,y)
  \label{eq:log_decomposition}
\end{equation}

Karena penjumlahan bersifat linear terhadap Transformasi Fourier:

\begin{equation}
  Z(u,v) = \mathcal{F}_i(u,v) + \mathcal{F}_r(u,v)
  \label{eq:freq_decomposition}
\end{equation}

di mana $\mathcal{F}_i$ dan $\mathcal{F}_r$ adalah DFT dari $\ln i$ dan $\ln r$.

\subsection{Filter Homomorphic}

Filter yang digunakan berfungsi menekan frekuensi rendah (iluminasi) dan
meningkatkan frekuensi tinggi (reflektansi):

\begin{equation}
  H(u,v) = \bigl(\gamma_H - \gamma_L\bigr)
            \left[1 - e^{-c\,D^2(u,v)/D_0^2}\right] + \gamma_L
  \label{eq:homo_filter_func}
\end{equation}

di mana $\gamma_L < 1 < \gamma_H$, dan $c$ mengontrol kecuraman transisi.

\subsection{Alur Kerja Lengkap Homomorphic Filtering}

\begin{enumerate}
  \item Hitung logaritma: $z(x,y) = \ln[f(x,y)]$
        (Persamaan~\eqref{eq:log_decomposition}).
  \item Hitung DFT: $Z(u,v) = \mathcal{F}\{z(x,y)\}$.
  \item Terapkan filter:
        \begin{equation}
          S(u,v) = H(u,v) \cdot Z(u,v)
          \label{eq:homo_apply}
        \end{equation}
  \item Hitung IDFT: $s(x,y) = \mathcal{F}^{-1}\{S(u,v)\}$.
  \item Kembalikan ke domain asli dengan eksponensial:
        \begin{equation}
          g(x,y) = e^{s(x,y)}
          \label{eq:homo_exp}
        \end{equation}
\end{enumerate}

\begin{tipbox}
\textbf{Kapan Menggunakan Homomorphic Filtering?}
\begin{itemize}
  \item Foto dengan pencahayaan tidak merata (bayangan satu sisi).
  \item Citra medis (X-ray, MRI) dengan rentang dinamis tinggi.
  \item Foto malam hari atau kondisi cahaya rendah.
\end{itemize}
Hasilnya: iluminasi dinormalisasi, detail tekstur (reflektansi) diperkuat.
\end{tipbox}

\subsection{Contoh Kode: Homomorphic Filtering}

\begin{lstlisting}[language=Python, caption={Implementasi lengkap homomorphic filtering}]
import numpy as np
from PIL import Image
import matplotlib.pyplot as plt

def homomorphic_filter(img, D0=30, gamma_L=0.5, gamma_H=2.0, c=1.0):
    """
    Homomorphic filtering untuk penyeimbangan iluminasi.
    img     : numpy array 2D, nilai [0, 255]
    D0      : cutoff frequency
    gamma_L : penekanan frekuensi rendah (< 1)
    gamma_H : penguatan frekuensi tinggi (> 1)
    c       : kecuraman transisi filter
    """
    M, N = img.shape

    # Langkah 1: Ln -- tambah 1 untuk menghindari log(0)
    ln_img = np.log1p(img.astype(np.float64))

    # Langkah 2: Centering + DFT
    x, y = np.arange(M), np.arange(N)
    xx, yy = np.meshgrid(x, y, indexing='ij')
    center_mask = (-1.0) ** (xx + yy)
    F = np.fft.fft2(ln_img * center_mask)

    # Langkah 3: Buat filter H(u,v) -- Persamaan homo_filter_func
    u = np.fft.fftfreq(M) * M
    v = np.fft.fftfreq(N) * N
    uu, vv = np.meshgrid(u, v, indexing='ij')
    D2 = uu**2 + vv**2
    H = (gamma_H - gamma_L) * (1 - np.exp(-c * D2 / (D0**2))) + gamma_L

    # Langkah 4: Filter + IDFT
    S = H * F
    s = np.real(np.fft.ifft2(S)) * center_mask

    # Langkah 5: Eksponensial -- kebalikan log1p
    g = np.expm1(s)
    return np.clip(g, 0, 255)

img = np.array(Image.open("image.jpg").convert("L"), dtype=np.float64)
result = homomorphic_filter(img, D0=30, gamma_L=0.5, gamma_H=2.0, c=1.0)

fig, axes = plt.subplots(1, 2, figsize=(10, 5))
axes[0].imshow(img.astype(np.uint8), cmap='gray')
axes[0].set_title('Citra Asli')
axes[1].imshow(result.astype(np.uint8), cmap='gray')
axes[1].set_title('Setelah Homomorphic Filtering')
for ax in axes:
    ax.axis('off')
plt.tight_layout()
plt.show()
\end{lstlisting}

% =============================================================================
\section{Ringkasan Perbandingan Filter Frekuensi}
% =============================================================================

\begin{center}
\renewcommand{\arraystretch}{1.5}
\begin{tabular}{llllp{4cm}}
\toprule
\textbf{Filter} & \textbf{Tipe} & \textbf{Tujuan} & \textbf{Ringing} & \textbf{Catatan} \\
\midrule
ILPF        & Lowpass  & Smoothing         & Parah    & Cutoff sempurna \\
BLPF        & Lowpass  & Smoothing         & Sedang   & Dikontrol orde $n$ \\
GLPF        & Lowpass  & Smoothing         & Tidak ada & Direkomendasikan \\
IHPF        & Highpass & Sharpening        & Parah    & --- \\
BHPF        & Highpass & Sharpening        & Sedang   & --- \\
GHPF        & Highpass & Sharpening        & Tidak ada & Direkomendasikan \\
HFE         & Highpass & Sharpening        & Sedang   & Pertahankan tonalitas \\
Homomorphic & Campuran & Normalisasi iluminasi & Minimal & Perlu preprocessing log \\
\bottomrule
\end{tabular}
\end{center}

% =============================================================================
\section{Mini-Latihan (Active Learning)}
% =============================================================================

\begin{enumerate}[leftmargin=*, label=\textbf{Latihan \arabic*.}]

  \item \textbf{Perhitungan Padding Minimal} \\
  Jika Anda memiliki citra berukuran $512 \times 512$ dan ingin melakukan pemfilteran
  di domain frekuensi, berapa nilai minimal $P$ dan $Q$ berdasarkan
  Persamaan~\eqref{eq:padding_size} agar terhindar dari \textit{wraparound error}?
  Mengapa dalam praktiknya kita biasanya membulatkan ke pangkat 2 terdekat?

  \begin{tipbox}
  \textbf{Petunjuk:} Hitung $P \geq 2(512) - 1 = 1023$. Pangkat 2 terdekat
  di atasnya adalah $1024$. Bagaimana kaitannya dengan efisiensi algoritma
  FFT Cooley-Tukey?
  \end{tipbox}

  \item \textbf{Mengapa GLPF Bebas Ringing?} \\
  Jelaskan secara matematis mengapa GLPF (Persamaan~\eqref{eq:glpf}) tidak
  menghasilkan efek \textit{ringing}, sementara ILPF (Persamaan~\eqref{eq:ilpf})
  menghasilkannya. Gunakan Teorema Konvolusi (Persamaan~\eqref{eq:conv_theorem})
  dalam penjelasan Anda.

  \item \textbf{Eksperimen Cutoff Frequency} \\
  Dengan menggunakan kode di Bagian 6.5, coba nilai $D_0 \in \{10, 30, 80, 150\}$
  untuk GLPF. Dokumentasikan bagaimana $D_0$ memengaruhi: (a) tingkat kekaburan
  citra, (b) noise yang tersisa, dan (c) ketajaman tepi.

  \item \textbf{Fase vs.\ Amplitudo} \\
  Rekonstruksi sebuah citra menggunakan (a) hanya informasi fase dan (b) hanya
  amplitudo. Amati hasilnya dan jelaskan mengapa fase lebih penting dalam
  merepresentasikan struktur visual.

  \begin{tipbox}
  \textbf{Kode Bantuan:}
  \begin{lstlisting}[language=Python]
F = np.fft.fft2(img)
# (a) hanya fase: set amplitudo = 1
phase_only = np.real(np.fft.ifft2(np.exp(1j * np.angle(F))))
# (b) hanya amplitudo: set fase = 0
mag_only   = np.real(np.fft.ifft2(np.abs(F)))
  \end{lstlisting}
  \end{tipbox}

  \item \textbf{Aplikasi Medis: GLPF vs.\ ILPF} \\
  Jelaskan mengapa dalam aplikasi medis (citra MRI atau mammogram), GLPF
  lebih disukai daripada ILPF untuk \textit{smoothing}. Hubungkan jawaban
  Anda dengan fenomena \textit{ringing} dan risiko artefak yang bisa
  disalahartikan sebagai kelainan medis.

\end{enumerate}

\end{document}